\centerline{\bf \Huge Комби-ТЧ}

\begin{flushright}
        \scriptsize{Эпизод 11. В полной темноте}
\end{flushright}

\problem{1}
Какое наибольшее количество различных натуральных чисел можно выбрать из множества $\{1,2,3, \ldots, n\}$ так, чтобы никакие два выбранных числа не делились друг на друга?

\problem{2}
Пусть $a_1, a_2, \ldots, a_{10}$ - натуральные числа, $a_1<a_2<\ldots<a_{10}$. Пусть $b_k-$ наибольший делитель $a_k$, меньший $a_k$. Оказалось, что $b_1>b_2>\ldots>b_{10}$. Докажите, что $a_{10}>500$.

\problem{3}
Существует ли такое натуральное число $n>2$, что все $n$ натуральных чисел от 1 до $n$ можно расставить по кругу в каком-то порядке таким образом, что произведение любых двух соседних чисел, увеличенное на 1 , будет кубом натурального числа?

\problem{4}
Дано натуральное число $n<10^{22}$. Докажите, что число $n$ может быть представлено в виде произведения десяти натуральных чисел таких, что среди них нет ни одного составного числа, большего 10000.

\problem{5}
Дано натуральное $n$. Докажите, что число перестановок $a_1, a_2, \ldots$, $a_n$ чисел $1,2, \ldots, n$, для которых при всех $i=1,2, \ldots, n$ число $i+a_i$ - простое, является квадратом целого числа.

\problem{6}
Дано 1000-значное число без нулей в записи. Докажите, что из этого числа можно вычеркнуть несколько (возможно, ни одной) последних цифр так, чтобы получившееся число не было натуральной степенью числа, меньшего 500 .

\problem{7}
Дано натуральное число $n$, свободное от квадратов. Пусть $D$ множество всех натуральных делителей $n$. Подмножества $A, B$ множества $D$ удовлетворяют условию, что для любых $a \in A$ и $b \in B$ справедливо, что $a$ не делится на $b$ и $b$ не делится на $a$. Докажите, что $\sqrt{|A|}+\sqrt{|B|} \leqslant \sqrt{|D|}$.